\section{Routing choices dictated by the observer}
\label{ba:ordered-routes}

The preceding processors move every branch by a prescribed affine map.
We now ask which parts of that motion must be controlled between the
integer times.  A bounded observation box allows horizontal work above
the observer.  A slab observer requires a bound on the horizontal
coordinate throughout the motion.  Finally, a half-space observer can be
protected by moving all sources to storage before delivering any target.
These three requirements lead to the constructions below.

\subsection{Four motions above a bounded observer}
\label{ba:four-stage-section}

We first separate ascent, scaling, translation and descent.  The two
planar motions occur at a fixed height, so their exclusion from the
bounded observer follows without a horizontal coordinate estimate.

Let $D_i,E_i\subset\mathbb R^2$, $1\le i\le N$, be closed rational
axis-aligned rectangles with positive side lengths.  Assume the $D_i$
are pairwise disjoint and the $E_i$ are pairwise disjoint, and let
\[
 A_i(b)=c_i'+\operatorname{diag}(\lambda_i,\lambda_i^{-1})(b-c_i),
 \qquad \lambda_i\in\mathbb Q_{>0},\qquad A_i(D_i)=E_i,
\]
where $c_i,c_i'$ are their centers.  No separation between a source and
a target is assumed.  All distances used to choose coordinate padding
below are distances for the maximum norm.

Let $\vartheta=\sigma$ be the effective flat switch in
\eqref{eq:p2-step}.  For a closed interval, possibly reduced to one
point, define
\[
 \chi_{[a,b],\epsilon}(v)=
 \vartheta\!\left(\frac{2(v-a+\epsilon)}\epsilon\right)
 \vartheta\!\left(\frac{2(b+\epsilon-v)}\epsilon\right).
\]
For a box $K$, let $\chi_{K,\epsilon}$ be the product over its three
coordinates.  It is one on the $\epsilon/2$ enlargement of $K$ and
zero outside its $\epsilon$ enlargement.  The effective profile estimates
following \eqref{eq:p2-step} apply to these products and rescalings.

\begin{lemma}[Four-stage realization]\label{ba:four-stage}
The rectangle data determine effectively a smooth divergence-free field
$W(s,x)$ supported in a compact subset of $(0,1)\times\mathbb R^3$.
Its time-one map is $(b,z)\mapsto(A_i(b),z)$ on a neighborhood of
each $D_i\times\{0\}$.  For a point of that plate, every planar
motion occurs at height $4i$, while ascent and descent preserve its
source and target planar coordinates, respectively.  The scaling can
use either linear interpolation of its positive first factor or linear
interpolation of the logarithm of that factor.
\end{lemma}

\begin{proof}\label{ba:four-stage-proof}
The empty family is realized by zero.  Otherwise put $h_i=4i$,
$w_i=c_i'-c_i$, and
\[
 \Theta_k(s)=\vartheta(10s-(2k-1)),\qquad 1\le k\le4.
\]
Only the $k$th motion is active on the interval
$[(2k-1)/10,2k/10]$; the four intervals are disjoint.  If
$a_{ij},a'_{ij}$ are the source and target half-widths, respectively,
take the four swept boxes
\[
 \begin{aligned}
 K_{i1}&=D_i\times[0,h_i],\\
 K_{i2}&=\{c_i+(u_1,u_2):|u_j|\le\max(a_{ij},a'_{ij})\}
                                      \times\{h_i\},\\
 K_{i3}&=\operatorname{box}(E_i-w_i,E_i)\times\{h_i\},\\
 K_{i4}&=E_i\times[0,h_i].
 \end{aligned}
\]
Here $\operatorname{box}$ denotes the smallest axis-aligned rectangle
containing the two displayed rectangles.  Within each of the four
families these boxes are disjoint: source separation handles the first,
target separation the fourth, and distinct heights the other two.
Choose a positive rational $\epsilon<1$ smaller than one quarter of
every pairwise distance within these families.  A family with fewer than
two members contributes no restriction.  Its $\epsilon$ enlargements
are then still disjoint.

Choose either
\[
 \mu_i=1+(\lambda_i-1)\Theta_2
 \quad\hbox{or}\quad \mu_i=\lambda_i^{\Theta_2},
 \qquad \gamma_i=\dot\mu_i/\mu_i.
\]
The factors $\mu_i$ and $\mu_i^{-1}$ stay between their respective
endpoint values.  Use the coordinate potentials
\[
 \begin{aligned}
 B_{i1}&=(0,h_i x_1,0),&
 B_{i2}&=(0,0,(x_1-c_{i1})(x_2-c_{i2})),\\
 B_{i3}&=(0,0,w_{i1}x_2-w_{i2}x_1),&
 B_{i4}&=(0,-h_i x_1,0).
 \end{aligned}
\]
For the logarithmic-scale construction we also retain the alternative
translation potentials, with $e_3=(0,0,1)$,
\[
 B_{i1}=\tfrac12(h_ie_3)\times x,\qquad
 B_{i3}=\tfrac12(w_{i1},w_{i2},0)\times x,\qquad
 B_{i4}=\tfrac12(-h_ie_3)\times x.
\]
Their curls equal the same constant translation vectors.  Their localized
fields need not agree away from the plates.  Either choice gives
\[
 W=\sum_i\left[
 \dot\Theta_1\nabla\times(\chi_{K_{i1},\epsilon}B_{i1})
 +\gamma_i\nabla\times(\chi_{K_{i2},\epsilon}B_{i2})
 +\dot\Theta_3\nabla\times(\chi_{K_{i3},\epsilon}B_{i3})
 +\dot\Theta_4\nabla\times(\chi_{K_{i4},\epsilon}B_{i4})\right].
\]
This is a finite sum of spatial curls.  Its spatial support is compact,
its time support is contained in $[1/10,8/10]$, and the lower bound
$\mu_i\ge\min(1,\lambda_i)>0$ proves smoothness and effectivity.

For $b\in D_i$, the successive paths are
\[
 \begin{array}{c|c}
 \text{motion}&\text{position}\\\hline
 \text{ascent}&(b,h_i\Theta_1)\\
 \text{scaling}&(c_i+\operatorname{diag}(\mu_i,\mu_i^{-1})(b-c_i),h_i)\\
 \text{translation}&(A_i(b)-w_i+\Theta_3w_i,h_i)\\
 \text{descent}&(A_i(b),h_i(1-\Theta_4)).
 \end{array}
\]
Each path stays in its swept box.  On a neighborhood of that box the
uncut curls are, in order,
\[
 h_ie_3,\qquad (x_1-c_{i1},-(x_2-c_{i2}),0),\qquad
 (w_{i1},w_{i2},0),\qquad -h_ie_3.
\]
After multiplication by the time coefficients, these are the prescribed
affine velocities.  The motion matrices are positive diagonal of determinant
one, the plateau margin is $\epsilon/2$, and the support padding is
$\epsilon$.  The disjoint swept boxes therefore permit the plateau
argument of Lemma~\ref{lem:bb-curl}, with $\eta=\epsilon/2$.
It identifies these particular localized fields with the displayed
paths on the plateaux.  The normalized diagonal factors are bounded by
$L_i=\max(1,\lambda_i,\lambda_i^{-1})$, so that lemma gives the
full initial neighborhood of radius $\epsilon/(4L_i)$.
\end{proof}

We use the following precise Euclidean comparison class in the applications:
on each $[0,T]$,
\begin{equation}\label{ba:euclidean-comparison-class}
 \begin{gathered}
 u\in C([0,T];H^2(\mathbb R^3))\cap C^1([0,T];L^2(\mathbb R^3)),\\
 u,\nabla u\in L^\infty([0,T]\times\mathbb R^3),\qquad
 p\in C([0,T];H^1(\mathbb R^3)).
 \end{gathered}
\end{equation}
The pressure representative is part of this condition; it is zero for
our constructed solution.  These comparison velocities and pressures
satisfy case (i) of Lemma~\ref{lem:b-comparison}.  We retain the stated
bounded-gradient requirement here, although that case does not need it.

That lemma also gives a smooth material diffeomorphism on every finite
interval.  Differentiating its trajectory equation gives
$\dot J=(D_xu)(t,\Phi_t(a))J$, $J(0)=I$, for $J=D_a\Phi_t$.
Consequently
\begin{equation}\label{ba:material-volume}
 \frac d{dt}\det J=(\nabla\cdot u)(t,\Phi_t(a))\det J=0,
 \qquad \det J=1.
\end{equation}
The same proof applies to periodic lifts on the torus.  The residual
formulas below are affine in viscosity, so they also define the asserted
solutions for any real $\nu>0$, with effectivity relative to $\nu$.

\begin{theorem}[Two bounded-box material tests]\label{ba:bounded-box-tests}
Fix a positive computable viscosity $\nu$.  A deterministic one-tape
machine and finite input effectively determine each of the following
general body forces on $[0,\infty)\times\mathbb R^3$.  Both have zero
initial velocity, a unique global smooth velocity in
\eqref{ba:euclidean-comparison-class}, and constructed pressure zero.
\begin{enumerate}
\item The force and velocity have one fixed compact spatial support, all
mixed derivatives bounded, and period one for $t\ge1$.  The particle
from $0$ enters $(-5,-2)\times(-1,2)\times(-1,1)$ exactly when the
machine halts.
\item The force and velocity have one fixed compact spatial support and,
for all $k\ge0$ and spatial multi-indices $\alpha$, satisfy
\begin{equation}\label{ba:eighteen-decay}
 \|\partial_t^k\partial_x^\alpha u(t)\|_\infty+
 \|\partial_t^k\partial_x^\alpha f(t)\|_\infty
 \le C_{k,\alpha}(1+t)^{-1-k}.
\end{equation}
The particle from $(2,0,0)$ enters
$(-1/2,3/2)\times(-1/2,3/2)\times(-1/2,1/2)$ exactly when the
machine halts.  In particular $f\in L^2([0,\infty);H^m(\mathbb R^3))$
for every integer $m\ge0$.
\end{enumerate}
The supports and constants may depend on the machine and input.  The
labels and observers do not.  All force derivatives admit evaluation to
any prescribed rational error by a finite program.
\end{theorem}

\begin{proof}\label{ba:bounded-box-tests-proof}
Normalize to a single halt state $q_H$ and a complete nonhalting table,
as in Section~\ref{ba:codes}.  Apply the move-first recorder of
Lemma~\ref{lem:bb-movefirst}.  Its history frontier starts at cell $2$;
all other history cells are empty and all mark bits are zero.  At
checkpoint $n$, the frontier is $g_n=n+2$, the work head satisfies
$|h_n|\le n$, and the next work instruction ending at $h'$ is completed
in exactly $4+2(g_n-h')$ local steps.  Thus every required checkpoint
is reached after finitely many steps.  The incoming-direction and
written-symbol inverse properties hold on the whole partial table.
The records recover the successive work instructions and hence the
absolute head position, even though the geometric code is head-relative.

Let $m$ be the full alphabet size.  Use odd digits $\beta(a)$ in
$\{1,3,\ldots,2m-1\}$, with base $B=2m+1$ in the first construction
and $B=2m+2$ in the second.  Set
\[
 C(a_0a_1\cdots)=\sum_{r\ge0}\beta(a_r)B^{-r-1},\quad
 P_a(v)=\frac{\beta(a)+v}{B},\quad D_a(v)=Bv-\beta(a),\quad
 I_a=P_a([0,1]).
\]
Let $\xi$ encode the left string, nearest cell first, and $\eta$
the string starting at the head.  The complete right, stay and left
branch maps are \eqref{eq:bb-one-cell-branches}, with its digit
$d(a)$ equal to $\beta(a)$ and its move $e$ equal to $d$ here.
The left move is split over every full symbol $c$.  The shared lemma
proves the full-rectangle endpoint formulas, arbitrary-tail action and
linear part $\operatorname{diag}(B^{-d},B^d)$.

Translate each state square horizontally.  In the first construction
give the halt state offset $-4$ and the other states offsets $2+2r$,
$r\ge0$; in the second use halt offset $0$ and other offsets $3r$,
$r\ge1$.  Thus the code is $(o_s+\xi,\eta,0)$.
Lemma~\ref{lem:bb-onecell} applies with exactly these odd digits and
either displayed base: its incoming conditions are those supplied by
the move-first recorder, and \eqref{eq:bb-one-cell-branches} is its minimal
left-split branch list.  Thus both complete closed rectangle families
are separately disjoint, with within-state gaps at least $B^{-2}$.
This proves the hypotheses of Lemma~\ref{ba:four-stage}; no separation
between the source and target families is required.

Use its linear scale and coordinate potentials for the first construction;
use its logarithmic scale and symmetric translation potentials for the
second.  In the latter case the scaling coefficient is exactly
$(\log\lambda_i)\dot\Theta_2$ and the scale matrix is
$\operatorname{diag}(\lambda_i^{\Theta_2},\lambda_i^{-\Theta_2})$.
Denote the resulting step field by $W_{\rm step}$.

The initial code $P_{\rm in}$ is rational.  A constant full-symbol
tail of digit $\beta_0$ after $M$ positions contributes
$B^{-M}\beta_0/(B-1)$; all other contributions are finite.  For the
appropriate fixed label $a_*$, put $v_0=P_{\rm in}-a_*$ and let
$J_0$ be the smallest box containing $[a_*,P_{\rm in}]$.  The loader
\[
 W_{\rm in}(s,x)=\frac d{ds}\vartheta(2s-1/2)\,
 \nabla\times\left(\chi_{J_0,1}(x)\frac{v_0\times x}{2}\right)
\]
has the exact marked path $a_*+\vartheta(2s-1/2)v_0$, since the
uncut curl is $v_0$ and the entire segment lies in the plateau.  Form
\[
 W(s,x)=W_{\rm in}(s,x)+\sum_{n\ge1}W_{\rm step}(s-n,x).
\]
The zero time collars make this smooth, initially zero, and periodic for
$s\ge1$, with common compact spatial support and bounded mixed derivatives.

At time $s=1+k$ the marked particle is exactly the code after $k$ local
instructions, through the first halt or indefinitely in a nonhalting run.
This follows by induction from the full-rectangle maps.  Each halt code
lies in its stated observation box, including one reached by loading an
initially halted machine.  In a nonhalting run, the first loader has
$x_1\ge0$ and later nonhalt codes have $x_1\ge2$; for the second these
bounds are $x_1\ge2$ and $x_1\ge3$.  Ascent and descent leave those
planar coordinates unchanged.  Every planar motion has height $4i\ge4$,
outside both observers.  The intervening stationary intervals add no new
positions.  This proves exclusion at every real time.

For the first case take $u=W$ and its residual force.  For the second set
\[
 s(t)=\log(1+t),\quad \theta(t)=(1+t)^{-1},\quad
 u(t,x)=\theta(t)W(s(t),x),
\]
whose force is
\[
 f=\theta^2\{-W+\partial_sW+(W\cdot\nabla)W\}(s(t),x)
                         -\nu\theta\Delta W(s(t),x).
\]
Lemma~\ref{ba:log-clock} applies to this fixed-support, bounded-derivative
template.  Its identity
$\partial_t(\theta^jY(s(t)))=\theta^{j+1}(\partial_s-j)Y(s(t))$
proves \eqref{ba:eighteen-decay} for every $k,\alpha$.  Fixed support
then gives the stated $L^2_tH^m_x$ conclusion.  Its clock is onto, with
inverse $t=e^s-1$, so the material path is the old path composed with
$s(t)$ and the event is unchanged.

In both cases smoothness, fixed support and bounded derivatives on
finite time intervals put the reference velocity in $CH^2\cap C^1L^2$
with bounded velocity and gradient.  It is solenoidal and initially zero.
Lemma~\ref{lem:b-comparison}(i) therefore gives the exact solution
$(u,0)$ and uniqueness in \eqref{ba:euclidean-comparison-class}.
The effective profiles following \eqref{eq:p2-step}, applied to the
finite table, rational initial code, finitely many boxes and repeated
template, give all derivative evaluations.  A coarse time approximation
encloses finitely many possible active translates, avoiding equality
tests at seams.  The prescription never computes the machine's run.
\end{proof}

\subsection{Small plates and a slab observer on the unit torus}
\label{ba:torus-plates-section}

A slab observer has no height restriction, so elevation alone cannot
protect it.  We now keep each nonhalting branch inside a narrow range of
first coordinates throughout its motion.  The delayed-head recorder also
allows arbitrary finitely supported work data on both sides of the head.

\begin{theorem}[Torus plates and square-integrable forcing]
\label{ba:torus-plates}
Fix a positive computable $\nu$.  Every deterministic one-tape machine
with finitely specified tape, blank outside a finite set, effectively
determines a smooth mean-zero general force on
$\mathbb T^3=\mathbb R^3/\mathbb Z^3$, with bounded mixed derivatives
and period one for $t\ge1$.  Its zero-data Navier--Stokes solution is
globally smooth and unique among smooth solutions, with pressure normalized
to mean zero and constructed pressure zero.  For
$P_*=(1/4,1/2,1/4)$ and $X(t)=\Phi_t(P_*)$, the event
\begin{equation}\label{ba:torus-plate-event}
 \exists t\ge0:\quad X_1(t)\pmod1\in(1/2,7/8)
\end{equation}
is equivalent to halting.

A separate effective force has the same zero data, event, uniqueness and
mean-zero property, is supported in one compact sub-box of $(0,1)^3$,
has all mixed derivatives bounded, and satisfies
\begin{equation}\label{ba:cube-rate}
 \|\widehat f(t,\cdot)\|_\infty=O((1+t)^{-2/3}).
\end{equation}
In particular $\widehat f\in L^2([0,\infty)\times\mathbb T^3)$.
\end{theorem}

\begin{proof}\label{ba:torus-plates-proof}
Use Lemma~\ref{ba:six-compiler} with the frontier initially at cell $1$.
Its main control $M_{q,e}$ records the last work displacement $e$, initially
zero.  In one work step it writes and marks the old head, records
$(q,e,a)$ at the frontier, returns to the mark, and only then makes the
work displacement.  All other initial history cells are blank and every
mark is zero.  At checkpoint $n$ the frontier is $n+1$, $|h_n|\le n$,
and the next checkpoint takes exactly $2(n+1-h_n)+3$ local steps.
These scans test history and marks, so the proof is unaffected by work
symbols on either side of the head.  Main controls with halting $q$ have
no outgoing rules.  The full-domain incoming properties are those of
the cited lemma.

Let $S$ be the local state set, $N=|S|$, let $S_h$ be its halting main
states, and enumerate $S$ by $n(s)=0,\ldots,N-1$.  Put
\[
 \ell=\frac1{16(N+1)},\quad z_0=\frac14,\quad
 x_s=\frac18+2\ell n(s)+\frac12\mathbf1_{s\in S_h},\quad
 T_s(\xi,\eta)=(x_s+\ell\xi,1/4+\ell\eta).
\]
Nonhalting squares lie in $1/8\le x_1<1/4$ and halting squares in
$5/8\le x_1<3/4$; distinct squares have gap at least $\ell$.
Give the full alphabet, of size $m$, odd digits in base $\beta=2m+1$.
Conjugating \eqref{eq:bb-one-cell-branches} by the source and target charts
gives separately disjoint rectangle families $R_j^0,R_j^1$ with positive
diagonal factors $\alpha_j,\alpha_j^{-1}$,
$\alpha_j\in\{\beta^{-1},1,\beta\}$.  Every chart has the same
scale, so this remains the physical linear part.
Lemma~\ref{lem:bb-onecell}, with this uniform chart scaling, gives
within-family coordinate gaps at least $\ell/\beta^2$, including the
two-letter target intervals for left moves.

For $1\le j\le L$, write the centers as $p_j^0,p_j^1$, the source
half-widths as $a_j,b_j$, and the target half-widths as
$\alpha_ja_j,\alpha_j^{-1}b_j$.  All endpoint half-widths are at
most $\ell/2$.  Choose
\[
 h_j=\frac12+\frac{j}{4(L+1)},\qquad
 \epsilon=\frac1{100}\min\left(1,\frac\ell{\beta^2},
                                      \frac1{4(L+1)}\right).
\]
With $\theta(s)=\vartheta(2s-1/2)$ and
$\Theta_k(\tau)=\theta(3\tau-k+1)$ for $k=1,2,3$, define
\[
 \begin{aligned}
 C_j(\tau)&=\big((1-\Theta_2)p_j^0+\Theta_2p_j^1,
                      z_0+(\Theta_1-\Theta_3)(h_j-z_0)\big),\\
 \sigma_j&=1+(\alpha_j-1)\Theta_2,\qquad
 \lambda_j=\dot\sigma_j/\sigma_j,\\
 (w_{j1},w_{j2},w_{j3})&=(a_j\sigma_j,b_j/\sigma_j,0).
 \end{aligned}
\]
These centers and half-widths describe moving plates.  Their
$\epsilon$ enlargements are disjoint: the first third uses source
projections, the middle third distinct heights, and the last third
target projections.  The padding is at most one hundredth of every
relevant gap.  All plates and padding lie in a compact sub-box of
$(0,1)^3$, since their centers interpolate between the stated chart
centers, widths are at most $\ell$, heights lie between $1/4$ and
$3/4$, and $\ell\le1/32$, $\epsilon\le1/100$.

For $w\ge0$ put
\[
 \kappa_\epsilon(w,v)=
 \vartheta\!\left(\frac{2(w+\epsilon-v)}\epsilon\right)
 \vartheta\!\left(\frac{2(w+\epsilon+v)}\epsilon\right),\qquad
 \chi_j=\prod_{a=1}^3\kappa_\epsilon(w_{ja},x_a-C_{ja}).
\]
Writing $y=x-C_j$, use
\[
 A_j=\tfrac12\dot C_j\times y+(0,0,\lambda_jy_1y_2),
 \qquad V(\tau,x)=\sum_j\nabla\times(\chi_jA_j).
\]
Extend the potentials by zero outside the chart.  The support remains
away from its faces, so this gives a smooth torus field, zero near the
time endpoints.  On the whole plateau of plate $j$ its velocity is
$\dot C_j+(\lambda_jy_1,-\lambda_jy_2,0)$; other summands vanish
there.  Hence the path from $(p_j^0+(r,s),z_0)$ is exactly
\begin{equation}\label{ba:plate-path}
 C_j(\tau)+(\sigma_j(\tau)r,\sigma_j(\tau)^{-1}s,0),
 \qquad |r|\le a_j,\quad |s|\le b_j.
\end{equation}
It ends at the prescribed affine image.  The plateau argument of
Lemma~\ref{lem:bb-curl} applies with
$D_j=\operatorname{diag}(\sigma_j,\sigma_j^{-1},1)$,
$\eta=\epsilon/2$, and diagonal bound
$\max(1,\alpha_j,\alpha_j^{-1})$: the required enlarged-box separation
was just proved.  It gives the same motion on an initial neighborhood
of radius $\epsilon/[4\max(1,\alpha_j,\alpha_j^{-1})]$, with constant
normal displacement, and therefore includes boundary codes.

The initial code $P_{\rm in}$ is rational, since relative to constant
blank tails only finitely many work and frontier digits differ.  Load it
from $P_*$ along $C_*=(1-\theta)P_*+\theta P_{\rm in}$ using
\[
 V_0(t,x)=\nabla\times\left[
 \prod_{a=1}^3\kappa_{1/64}(0,x_a-C_{*,a}(t))
                \frac{\dot C_*(t)\times(x-C_*(t))}{2}\right].
\]
The segment and its collar stay within the chart: its first coordinate
stays between $1/8$ and $3/4$, its second between $1/4$ and
$1/2+\ell$, and its third is $1/4$.  Its plateau curl is $\dot C_*$.
Take $U=V_0$ on $[0,1]$ and $U(t)=V(t-n)$ on $[n,n+1]$ for $n\ge1$.
Time collars make this smooth and initially zero.  It has fixed compact
chart support, bounded mixed derivatives, and period one after loading.
Define $f=U_t+(U\cdot\nabla)U-\nu\Delta U$ using the stated physical
viscosity.  Both $U$ and $f$ have mean zero: $U$ is a curl, the
Laplacian integrates to zero, and
$(U\cdot\nabla)U=\operatorname{div}(U\otimes U)$.
The torus clause of Lemma~\ref{lem:b-comparison} applies to this
smooth, initially zero, solenoidal reference field and gives the unique
smooth solution $(U,0)$ with mean-zero pressure.

Equations~\eqref{eq:bb-one-cell-branches} and \eqref{ba:plate-path} identify
the particle at time $1+k$ with the $k$th local configuration.  The
recorder invariant ensures that a nonhalting computation continues
indefinitely.  A halt, including an initial halt reached by loading,
puts the particle in $5/8\le X_1<3/4$, proving the event.  For a
nonhalting computation the loader has $X_1\le1/4$ and every later
branch has both endpoint centers below $1/4$.  Its center is their
convex interpolation and its first half-width is at most $\ell/2$.
Thus $0<X_1<1/4+\ell/2<1/2$ at every intermediate time.  The path
stays inside the chart, so reduction modulo one cannot create a false visit.

For the separate decaying force choose the onto clock
\[
 s(t)=(1+t)^{1/3}-1,\qquad \widehat U(t,x)=s'(t)U(s(t),x).
\]
Its residual is
\[
 \widehat f=s''U+(s')^2\{\partial_sU+(U\cdot\nabla)U\}
                                         -\nu s'\Delta U,
\]
where the fields on the right are evaluated at $(s(t),x)$.  Since
$s'=(1/3)(1+t)^{-2/3}$ and $s''=-(2/9)(1+t)^{-5/3}$, the bounded
template derivatives prove \eqref{ba:cube-rate}.  Every positive-order
derivative of $s$ is bounded; repeated product and chain rules therefore
bound every mixed derivative of $\widehat U$ and $\widehat f$.  The
square of $(1+t)^{-2/3}$ is integrable.  Support and mean zero persist,
and $U(0)=0$ retains the zero datum.  The torus clause of
Lemma~\ref{lem:b-comparison} therefore applies to $\widehat U$ too.
Finally the material path is the old one at $s(t)$,
and $s$ has inverse $t=(1+s)^3-1$, so the event is unchanged.

All choices are effective finite data followed by explicit smooth
formulas.  In particular $\sigma_j\ge\min(1,\alpha_j)>0$; the cutoff
estimates and zero collars allow derivative evaluation at approximate real
arguments without tests for exact boundary equality.  Only positive
arguments enter the cube root.  The program repeats the entire local
table and never uses a computed execution.
\end{proof}

\subsection{Storage and periodic forcing without a loader}
\label{ba:parking-section}

The next construction absorbs initialization into the positions of the
state rectangles.  A zero digit for the full blank symbol places the
initial left stack at zero; an input-dependent second-coordinate translation then
places the entire initial code at one fixed particle.  To protect a
half-space observer during every period, we first evacuate all sources
to separate storage rectangles and only afterwards fill the targets.

\begin{lemma}[Zero-blank coding]\label{ba:zero-blank-code}
Let an alphabet of size $m>1$ have digits $0,\ldots,m-1$, with full
blank digit zero.  Set $B=m+1$, $\kappa=(m-1)/m$, and
\[
 \operatorname{val}(a_0a_1\cdots)=\sum_{r\ge0}a_rB^{-r-1},\qquad
 I_a=\frac{a+[0,\kappa]}B,\qquad I_{a,b}=\frac{a+I_b}B.
\]
The code is injective, takes values in $[0,\kappa]$, and its distinct
one-letter and two-letter cylinders have positive rational gaps.  For a
deterministic partial table with target-state-determined incoming direction
and incoming rule determined by target state and written symbol, separated
state translates of $[0,\kappa]^2$ yield separately disjoint full source
and target rectangles by the following maps in state-relative coordinates:
\begin{equation}\label{ba:zero-digit-maps}
 \begin{array}{c|c|c|c}
 d&\text{source}&(\xi',\eta')&\text{target}\\\hline
 1&[0,\kappa]\times I_a&((b+\xi)/B,B\eta-a)&I_b\times[0,\kappa]\\
 0&[0,\kappa]\times I_a&(\xi,\eta+(b-a)/B)&[0,\kappa]\times I_b\\
 -1&I_c\times I_a&(B\xi-c,c/B+(b+B\eta-a)/B^2)&
                                      [0,\kappa]\times I_{c,b}.
 \end{array}
\end{equation}
The last row is split over all $c$.  Each map has positive diagonal
linear part of determinant one and acts exactly on the full rectangle.
\end{lemma}

\begin{proof}\label{ba:zero-blank-code-proof}
The maximum tail is $(m-1)/(B-1)=\kappa$, and
$(m-1+\kappa)/B=\kappa$.  Consecutive first-digit intervals have
gap $(1-\kappa)/B>0$.  Within one first-digit interval the second-digit
gaps are this quantity divided by $B$; different first digits already
give a gap.  The interval containing a code recovers its first digit,
and $v\mapsto Bv-a$ then exposes the tail.  Repetition proves
injectivity, including endpoint codes.

Writing $a$ to $b$ changes the first right digit.  A right move pushes
$b$ onto the left string and removes $a$ from the right; a left move
removes $c$ from the left and prepends $c,b$ to the old right tail.
These identities give \eqref{ba:zero-digit-maps}, including its images
and reciprocal slopes.  Determinism and digit gaps separate sources by
state, read digit, and $c$ when needed.  At one target state the common
incoming direction determines where the written digit is visible: in
$I_b$ on the left or right, or in $I_{c,b}$.  The incoming-rule
hypothesis and digit gaps then separate the closed target rectangles.
\end{proof}

\begin{lemma}[Evacuation, storage and delivery]\label{ba:parking}
Suppose $A_i:D_i\to E_i$, $1\le i\le N$, are the rational reciprocal
rectangle maps of Lemma~\ref{ba:four-stage}, with $N\ge1$ and every
$D_i\subset\{x_1<0\}$.  There is an effective smooth divergence-free
field on $\mathbb R\times\mathbb R^3$, one-periodic in time, zero near
each integer time, and uniformly compactly supported in space, whose
one-period map is $(b,z)\mapsto(A_i(b),z)$ on a neighborhood of every
$D_i\times\{0\}$.  If $E_i\subset\{x_1<0\}$, the entire
one-period path of each point in $D_i\times\{0\}$ remains in that
half-space.
\end{lemma}

\begin{proof}\label{ba:parking-proof}
Let $d_i,b_i$ be the source and target centers, and let $a_{ij},a'_{ij}$
be their half-widths.  Put $a^*_{ij}=\max(a_{ij},a'_{ij})$.
Choose $\epsilon$ to be one quarter of the minimum of $1$ and all
pairwise maximum-norm distances within the source and target families;
omit pair entries when $N=1$.  Define
\[
 \begin{gathered}
 x_{\min}=\min\{x_1:(x_1,x_2)\in D_i\cup E_i\text{ for some }i\},\\
 R=\max(1,\max_i a^*_{i1}),\qquad J=4(R+1),\qquad
 s_i=(x_{\min}-iJ,0).
 \end{gathered}
\]
The storage rectangle at $s_i$ with half-widths $a^*_{ij}$ lies in
$x_1<0$.  Its first-coordinate projection is separated by more than one
from every source and target and from the other storage projections:
the relevant gaps are at least $J-R>1$ and $J-2R>1$.

We need only a localized translation and a localized reciprocal scaling.
On a time segment $[t_0,t_0+\delta]$, use
$g(t)=\vartheta(3(t-t_0)/\delta-1)$, whose derivative is supported
in the middle third.  Using the cutoff $\kappa_\epsilon$ defined in
the proof of Theorem~\ref{ba:torus-plates}, for half-widths $a,b$ write
\[
 \chi_{a,b}(y)=\kappa_\epsilon(a,y_1)
               \kappa_\epsilon(b,y_2)\kappa_\epsilon(0,y_3).
\]
A plate translation from center $P$ to $Q$ is produced by
\[
 C(t)=(1-g)P+gQ,\qquad
 \nabla\times\left[
 \tfrac12\chi_{a,b}(x-C(t))\dot C(t)\times(x-C(t))\right].
\]
On its plateau the velocity is $\dot C$ and the exact path is
$C(t)+(r,s,0)$.  At storage center $P=(s_i,0)$, take
$q(t)=1+(\lambda_i-1)g(t)$ and the field
\[
 \nabla\times\left(0,0,
 \frac{\dot q}{q}(x_1-P_1)(x_2-P_2)
                           \chi_{a^*_{i1},a^*_{i2}}(x-P)\right).
\]
It has the exact path $P+(qr,q^{-1}s,0)$.  The denominator has the
positive lower bound $\min(1,\lambda_i)$, and the whole path lies in
the storage rectangle, because each half-width stays between its source
and target values.

Partition a period into $7N$ equal segments and take $Z=2$.  For each
$i$ in order, use four segments for
\[
 (d_i,0)\longrightarrow(d_i,Z)\longrightarrow(s_i,Z)
                \longrightarrow(s_i,0),\qquad\text{then scale at }s_i.
\]
Only after all these evacuations, use three segments for each $i$ to
deliver its target shape along
\[
 (s_i,0)\longrightarrow(s_i,Z)\longrightarrow(b_i,Z)
                                      \longrightarrow(b_i,0).
\]
Sum these segment fields and extend periodically.  The middle-third
time switches give zero neighborhoods of every segment endpoint and
every integer time.  The result is a smooth finite sum of curls with
one compact spatial support.

This is the $7N$ schedule of Lemma~\ref{lem:bb-lift-parking}(ii).
Its separation hypotheses have explicit margins here.  Unpadded
source-source and target-target gaps are at least $4\epsilon$, so their
$\epsilon$ enlargements leave gap at least $2\epsilon$.  Storage
gaps exceed one, leaving more than $1-2\epsilon>0$ after padding.
At horizontal travel height the enlarged moving plate and every
enlarged stationary plate have vertical gap $2-2\epsilon>0$.
Every scaling stays in its reserved storage rectangle.  The lemma
therefore gives simultaneous noninterference for evacuation and delivery,
including when a source meets a target.

The localized potentials above have exactly these affine generators on
their plateaux, so the plateau argument of Lemma~\ref{lem:bb-curl}
gives endpoint
$(d_i+(r,s),0)\mapsto(b_i+(\lambda_i r,\lambda_i^{-1}s),0)$,
with normal displacement unchanged.  Put
$M=\max_i(1,\lambda_i,\lambda_i^{-1})$ and
$d_*=\min(2\epsilon,1-2\epsilon,2-2\epsilon)>0$.
A neighborhood radius smaller than both $\epsilon/(8M)$ and
$d_*/(4M)$ retains the plateaux and all exclusion margins throughout
the concatenation.  This applies the shared proofs to the displayed
potentials and does not identify their fields away from the plateaux.

If $E_i$ is negative, each translation of an individual point joins two
negative first coordinates, and its linear interpolation remains negative.
Scaling occurs inside the negative storage rectangle.  The point is
stationary during all other segments.  This proves the asserted
whole-period exclusion.
\end{proof}

\begin{theorem}[A periodic force with no loading interval]
\label{ba:periodic-parking}
Fix a positive computable $\nu$.  A deterministic one-tape machine and
finite input word in cells $0,\ldots,L-1$, with other cells blank and
initial head zero, effectively determine a general body force on
$\mathbb R^3$ extending smoothly and one-periodically to all real times.
It has uniformly compact spatial support and bounded mixed derivatives.
Its zero-data solution has a unique global smooth velocity in
\eqref{ba:euclidean-comparison-class}, with constructed pressure zero, and
\[
 \exists t\ge0:\quad (\Phi_t((-2,0,0)))_1>0
 \quad\Longleftrightarrow\quad\text{the machine halts}.
\]
The fixed label is already the initial code; periodicity begins at zero.
\end{theorem}

\begin{proof}\label{ba:periodic-parking-proof}
Insert a fresh nonhalting start state that preserves the read letter,
stays put and enters the original start state.  Replace missing
nonhalting instructions by letter-preserving stay instructions to a halt
state.  This finite normalization preserves the halting question,
including an initially halted machine, and guarantees a nonhalting
initial control and at least one rule.

Apply the six-row delayed-head recorder of Lemma~\ref{ba:six-compiler}
with initial frontier $1$.  Initially every other history cell is blank
and every mark is zero.  The main control records $(q,e)$, where $e$
is the preceding displacement, and the history record is $(q,e,a)$.
The old head is marked before the right scan and cleared on return,
immediately before displacement into the next main state.  At checkpoint
$k$ the frontier is $k+1$ and the head $i$ satisfies $|i|\le k$;
the next checkpoint uses $2(k+1-i)+3\le4k+5$ local steps.  The
incoming-direction and written-symbol properties hold on the full
domain, as required by Lemma~\ref{ba:zero-blank-code}.

Give the full blank symbol digit zero and the other full symbols digits
$1,\ldots,m-1$.  Order nonhalting local states with the initial state
first and assign offsets $-2h$, $h=1,2,\ldots$; assign halting main
states offsets $2h$, $h=1,2,\ldots$.  Let $\eta_{\rm in}$ be the
initial right-stack value and put $y_*=-\eta_{\rm in}$.  Only finitely
many full symbols are initially nonblank, so this is a finite rational
sum.  Every cell left of the initial head is fully blank.  The code
\[
 (x_\sigma+\xi,y_*+\eta,0)
\]
therefore places the initial configuration exactly at $(-2,0,0)$.
The input changes the common second-coordinate translation, while the
observed label and zero initial velocity remain fixed.

Every finitely reached tape still has finite full-symbol support, so its
code is rational and can be decoded by rational digit extraction until
both tails are zero.  At a checkpoint the leftmost nonempty history site
has absolute index $1$; its relative position also recovers the head's
absolute index.  Thus this coding retains the initialized computation.

Translate the rectangles of \eqref{ba:zero-digit-maps} by the indicated
state offsets and $y_*$.  Lemma~\ref{ba:zero-blank-code} gives their
full affine action and separate source and target gaps.  Since
$\kappa<1$, every nonhalting state block is strictly negative in the
first coordinate and every halting block strictly positive.  Halting
states have no outgoing rules, so all sources are negative.  Apply
Lemma~\ref{ba:parking} to obtain its periodic field $v$.  Zero tails
can place a code on a rectangle boundary, which is included in that
lemma's full-neighborhood conclusion.

Set
\[
 f=f_0+\nu f_1,\qquad
 f_0=\partial_tv+(v\cdot\nabla)v,\qquad f_1=-\Delta v.
\]
The field and force are smooth, periodic, uniformly compactly supported,
and have all mixed derivatives bounded.  Since $v$ vanishes near integer
times, its initial value is zero.  Fixed support and bounded derivatives
give $v\in CH^2\cap C^1L^2$ and bounded $v,\nabla v$ on every finite
interval.  Lemma~\ref{lem:b-comparison}(i) therefore supplies the exact
solution $(v,0)$, uniqueness in
\eqref{ba:euclidean-comparison-class}, and its global material flow.

At time zero the marked particle is the initial code.  If it is a
nonhalting code at integer $n$, its digits select a source rectangle and
the next period applies exactly that local instruction.  Induction
identifies the integer-time codes through the first halt or forever.
Finite completion of every recorder scan identifies the original
machine at the checkpoint integers.  A halt therefore reaches a positive
code.  If the machine never halts, each used branch has a negative target
rectangle.  Lemma~\ref{ba:parking} keeps its marked path negative for
that whole period, ruling out an intermediate false visit.

Finally the force program uses finite normalization, the finite table,
the rational input shift, a finite rectangle list, rational storage and
separation data, $7N$ time segments and explicit cutoffs.  All scaling
denominators have known positive lower bounds.  The zero-collar and flat
profile estimates following \eqref{eq:p2-step} permit evaluation of the
repeated formula and every derivative from approximate real arguments
without exact seam tests.
Every period executes all branch maps, independently of the marked
particle's current state; the program never computes the input run.
\end{proof}

\subsection{Direction-split marking and a bounded observer}
\label{ba:bounded-box-native}

A stay instruction need not clear and reinstall a persistent mark.
We separate that case in the next table.  Its geometric realization
uses a bounded observation box: during horizontal transfers the routing
height itself prevents a false visit.

\begin{lemma}[Direction-split persistent marking]
\label{ba:split-mark-compiler}
In the notation of Lemma~\ref{ba:persistent-compiler}, there is an
effective partial finite table with a persistent checkpoint mark,
pointer at $p_n=n+2$, and both full-domain incoming properties of
Lemma~\ref{ba:rectangles}.  A work step ending at $w'$ takes
$2(p_n-w')+4+\mathbf1_{\{d_\tau\ne0\}}$ local steps.
\end{lemma}

\begin{proof}\label{ba:split-mark-compiler-proof}
Use full symbols $[g,e,l]$, with work letter $g$, mark $e\in\{0,1\}$,
and $l\in\{E,P\}\sqcup\{J_\tau\}$.  Controls are $A_q$,
$B_\tau$ for moving instructions, $C_\tau,S_\tau$, and $D_q,R_q$.
For $\tau=(q,a)$ writing $b_\tau$, the complete table is
\[
\begin{array}{c|c|c|c|r}
\text{control}&\text{read and guard}&\text{target}&\text{write}&d\\\hline
A_q&[a,1,l],\ d_\tau=0&C_\tau&[b_\tau,1,l]&0\\
A_q&[a,1,l],\ d_\tau\ne0&B_\tau&[b_\tau,0,l]&d_\tau\\
B_\tau&[g,0,l]&C_\tau&[g,1,l]&0\\
C_\tau&[g,1,l],\ l\ne P&S_\tau&[g,1,l]&1\\
S_\tau&[g,0,l],\ l\ne P&S_\tau&[g,0,l]&1\\
S_\tau&[g,0,P]&D_{q^+_\tau}&[g,0,J_\tau]&1\\
D_q&[g,0,E]&R_q&[g,0,P]&-1\\
R_q&[g,0,l],\ l\ne P&R_q&[g,0,l]&-1\\
R_q&[g,1,l],\ l\ne P&A_q&[g,1,l]&0.
\end{array}
\]
The first two rows use nonhalting $q$.  In those rows and the
$B_\tau$ row, $l$ is unrestricted, including $P$; no additional
history guard is imposed.  Initialize in $A_{q_0}$ at zero, with the
input work tape, a unique mark at zero, pointer at two, and empty
history elsewhere.

At checkpoint $n$, the work move ends at
$w'=w_n+d_\tau\le n+1<p_n$.  A stay instruction reaches $C_\tau$
in one step with the mark retained.  A moving instruction reaches it
in two steps after transferring the mark.  From $C_\tau$, move right,
scan the $p_n-w'-1$ intervening nonpointer cells, record at $p_n$,
and extend the pointer to $p_n+1$.  The return scans $p_n-w'$
unmarked cells and stops at the persistent mark, entering the correct
$A_{q^+_\tau}$.  There are five other steps for a stay instruction
and six for a moving instruction.  This yields the stated count and
restores the invariant.  Every required guard holds, so the execution
continues exactly until halting, including possible initial halting.

For an arbitrary allowed tape, each target determines its incoming move.
It is zero into $A_q$ or $C_\tau$, $d_\tau$ into $B_\tau$, one into
$S_\tau$ or $D_q$, and minus one into $R_q$.  The record $\tau$ recovers the
source of the work update.  A fixed $C_\tau$ receives only an $A_q$
entry if $d_\tau=0$, and only $B_\tau$ entries otherwise; the written
symbol recovers the old full symbol in either case.  At $S_\tau$,
the written mark separates entry from scan.  At $D_q$, $J_\tau$
identifies the recording instruction and its former pointer.  At $R_q$,
the written pointer separates entry from the loop.  At $A_q$, the
unique source is $R_q$ and the marked symbol is unchanged.  Undoing
the indicated shift therefore gives at most one predecessor in every
case, on the full stated domain.
\end{proof}

\begin{theorem}[A bounded-box event with a vertical routing guard]
\label{ba:bounded-box-euclidean}
For every machine and finite input, and fixed positive computable
viscosity, there is an effective smooth force on $\mathbb R^3$ with
uniformly compact spatial support, bounded mixed derivatives and period
one for $t\ge1$, whose zero-data solution satisfies
\[
 \exists t\ge0:\ X(t;(2,0,0))\in(0,1)\times(0,1)\times(-1,1)
 \quad\Longleftrightarrow\quad\text{the machine halts}.
\]
The constructed solution is smooth with pressure zero.  It is unique
among smooth solutions with $u\in C_tH^2\cap C_t^1L^2$, bounded
$u$ on each finite interval, and $p\in C_tH^1$.  This comparison class
requires no bound on the competing velocity's spatial gradient.
\end{theorem}

\begin{proof}\label{ba:bounded-box-euclidean-proof}
Put $\rho(s)=\sigma(2s-1/2)$, using the effective flat switch
from~\eqref{eq:p2-step}.
Use the single-halt normalization and
Lemma~\ref{ba:split-mark-compiler}.  For its full alphabet $\Sigma$,
choose distinct odd digits $e(a)\in\{1,3,\ldots,2|\Sigma|-1\}$ and
base $B=2|\Sigma|+1$.  For an infinite word
$\omega=\omega_0\omega_1\cdots$, define
\[
 c(\omega)=\sum_{j\ge0}e(\omega_j)B^{-j-1}.
\]
Place $A_h$ at offset zero and the
other control squares at first-coordinate offsets $3,6,9,\ldots$.
The code at height zero is
\[
 (o_s+c(\text{left stream}),\ c(\text{right stream}),0).
\]
Every tail value lies in
$[1/(B-1),(B-2)/(B-1)]\subset(0,1)$, so a halt code lies strictly
inside the observed box.  The two incoming properties give the
separated filled rectangles of Lemma~\ref{ba:rectangles}, with
minimal left subdivision and first factors $1,B^{-1},B$ for stay,
right and left moves.

Apply the linear version of Lemma~\ref{ba:two-interpolations} to
the full cuboids with third interval $[-1,1]$, heights $6i$, and
switches $\rho(3s),\rho(3s-1),\rho(3s-2)$.  Source projections
separate lifts, heights separate horizontal motion, and target
projections separate descents.  Equation~\eqref{ba:affine-box-field}
gives the exact affine flow on every swept cuboid and a neighborhood.

The initial code $x^{\rm in}$ is rational.  For example, if
$b=\max(3,\text{input length})$, the first $b$ right-stream symbols
include the entire input and the pointer at cell two, after which the
stream is constant; the left stream is constant.  On the loading
interval use
\[
 C_*(t)=(2,0,0)+\rho(t)(x^{\rm in}-(2,0,0))
\]
and
\[
 V_*(t,x)=\nabla\times\left\{
 \chi(C_*(t),(1,1,1),1;x)
 \tfrac12 C_*'(t)\times(x-C_*(t))\right\}.
\]
The cutoff is the centered one defined in
Lemma~\ref{ba:two-interpolations}.  Its plateau velocity is $C_*'$, so
the fixed particle follows the displayed path.  This field vanishes
outside $[1/4,3/4]$.  Extend the step field $V$ by zero and set
$U=V_*+\sum_{k\ge1}V(t-k)$, with its residual force at the prescribed
$\nu$.  The finite union of swept cuboids and the loading tube is
bounded; all time joins have zero collars.  Consequently $U$ and its
force are smooth, uniformly compactly supported, and have bounded
mixed derivatives, with period one after loading.  Effective evaluation
uses the finite formulas and a finite superset of potentially active
integer shifts at any approximate time, as in
Lemma~\ref{ba:realization}.

The precise uniqueness input is Lemma~\ref{lem:b-comparison}(i).
On each finite interval the constructed $U$ is smooth, has one compact
spatial support, and has bounded velocity and gradient.  Consequently
$U\in CH^2\cap C^1L^2$, $U(0)=0$, and the force is exactly
$\mathcal F_\nu[U]$.  The theorem's competitors are smooth classical
$v\in CH^2\cap C^1L^2$ with bounded $v$ and a representative
$p\in CH^1$, precisely case (i).  No bound on $\nabla v$ is added.
That case gives $v=U$ and fixes the $H^1$ pressure representative to
zero.  The lemma also supplies the unique global material trajectories
of the reference field.

Loading and the exact full-box maps put the particle at the code after
$k$ local steps at time $1+k$, through the first halt or forever in a
nonhalting execution.  A halt code, including an initially halted code
reached by loading, is inside the observation box.  In a nonhalting
run, loading has first coordinate at least two.  During a later lift,
the particle keeps its nonhalting source planar code, whose first
coordinate exceeds three.  During descent it keeps the nonhalting
target planar code with the same property.  Throughout the middle
third its height is $6i>1$.  These three exclusions cover every real
time and prove the event equivalence independently of horizontal
excursions in the middle stage.  The repeated processor installs every
local branch from the finite table; it does not depend on the executed
run or its outcome.
\end{proof}


\subsection{Separating vertical and planar potentials}

\begin{theorem}[A decaying processor with phasewise planar potentials]
\label{ba:phasewise-flow}
At fixed computable $\nu>0$ there is an effective force on $\mathbb R^3$
with one compact spatial support and all mixed derivatives of time order
$k$ bounded by $C_{k,\alpha}(1+t)^{-1-k}$, whose unique smooth
zero-data solution satisfies
\[
 \exists t\ge0:\ (\Phi_t(0))_1<-1/2
 \quad\Longleftrightarrow\quad\text{halting}.
\]
It uses the six-rule recorder with initial frontier 1, full subdivision
of the rectangles, and separate vertical-translation and planar
Hamiltonian potentials.  The constructed pressure is zero and uniqueness
is in Lemma~\ref{ba:realization}'s whole-space class.
\end{theorem}
\begin{proof}\label{ba:phasewise-flow-proof}
Let $\Gamma$ be the full alphabet of the six-rule table with $r_0=1$.
Use odd digits in base $B=2|\Gamma|+1$ and its fully subdivided maps.  Put halt lanes at $-2,-4,\ldots$
and the other lanes at $2,4,\ldots$, each of width one.  The second
coordinate lies in $[0,1]$.  For the resulting rectangles $P_i,Q_i$,
centers $p_i,q_i$, and first scale $\lambda_i$, choose rational
$0<\varepsilon\le1/4$ smaller than one quarter of all within-source
and within-target maximum-norm distances.  Empty pair lists impose no
condition.  Thicken by $[-1,1]$ and put $H_i=4i$.

On each third of the cycle use a flat progress function $\vartheta$
from zero to one.  Lifting uses $z_i=H_i\vartheta$ and potential
$A_i=(0,z_i'X_1,0)$, whose curl is $(0,0,z_i')$.
At height $H_i$, set $c_i=(1-\vartheta)p_i+\vartheta q_i$ and
$\mu_i=1-\vartheta+\vartheta\lambda_i$.  Use the potential
$A_i=(0,0,\psi_i)$, where
\[
 \psi_i=c_{i1}'(X_2-c_{i2})-c_{i2}'(X_1-c_{i1})
       +\frac{\mu_i'}{\mu_i}(X_1-c_{i1})(X_2-c_{i2}).
\]
Its curl is
$(c_{i1}'+(\mu_i'/\mu_i)(X_1-c_{i1}),
 c_{i2}'-(\mu_i'/\mu_i)(X_2-c_{i2}),0)$.
Thus it realizes reciprocal scaling about the interpolated center.
Descending uses the first potential with $z_i=H_i(1-\vartheta)$.
Multiply each potential by its current-box $\varepsilon$ cutoff
before taking the curl.  Source gaps, distinct middle heights, and
target gaps separate the three stages.  On the whole boxes and a
neighborhood the uncut formulas are exact.  The first coordinate is
constant during vertical motion and is the convex interpolation of its
endpoints during the planar stage.

For loading put $\theta_0(s)=\sigma(2s-1/2)$, which is zero near zero
and one near one. Let $P_*$ be the rational initial code. Load it in the
first unit interval by the planar path
$m(s)=\theta_0(s)((P_*)_1,(P_*)_2)$, localized around
$(m(s),0)+[-1,1]^3$, using potential
$(0,0,m_1'(X_2-m_2)-m_2'(X_1-m_1))$.
Repeat the three-phase instruction field afterward to obtain $W$.
Code induction and the delayed-head invariant prove the correspondence
at times $1+n$.  Nonhalting codes have first coordinate at least 2,
so loading and the convex-path bound exclude $X_1<-1/2$ for every real
time.  A halt code has first coordinate at most $-1$, including one
reached by loading an initially halted machine.  Apply
Lemma~\ref{ba:log-clock}; physical code times are $e^{1+n}-1$.
Its support, derivative, event, and uniqueness conclusions give precisely
the claimed separate force choice.
\end{proof}
